\editionPage{references}{Références officielles}{Recherche documentaire du 25 octobre 2026}{Sources de cadrage ; exercices et figures rédigés pour ce manuel.}
\editionEtape{Programme français}
Ministère de l'Éducation nationale, programme de physique-chimie de Seconde générale et technologique, arrêté du 17 janvier 2019, BO spécial n° 1 du 22 janvier 2019, NOR MENE1901634A. Annexe de 13 pages.\par\smallskip
\url{https://www.education.gouv.fr/bo/19/Special1/MENE1901634A.htm}
\editionEtape{Vérification du programme en vigueur}
Éduscol, « Programmes et ressources en physique-chimie - voie GT », page mise à jour en mars 2026, consultée le 25 octobre 2026. Le programme de Seconde y est référencé sous le BO de janvier 2019.\par\smallskip
\href{https://eduscol.education.gouv.fr/5829/programmes-et-ressources-en-physique-chimie-voie-gt}{Éduscol : programmes et ressources en physique-chimie, voie GT.}
\editionEtape{Ressources d'accompagnement consultées}
Ressources Éduscol sur le cortège électronique, le dénombrement des entités et les signaux électriques en Seconde ; fiches d'accompagnement pour Python : positions, vecteurs déplacement et vitesse, graphiques et variabilité des volumes. Elles éclairent le niveau attendu et la démarche de modélisation.
\editionEtape{Méthode éditoriale}
Le programme a été lu dans son intégralité pour établir la correspondance des notions et des capacités. Les activités du livre sont originales et leurs données sont adaptées à l'apprentissage. Les six devoirs ne sont pas des sujets nationaux : il s'agit d'évaluations de la collection, avec corrigés et barèmes proposés.

